\documentclass[a4paper,12pt]{article}

\usepackage[utf8]{inputenc}
\usepackage[T1]{fontenc}
\usepackage[table]{xcolor}
\usepackage{geometry}
\usepackage{longtable}
\usepackage{enumitem}
\usepackage{hyperref}
\usepackage{listings}
\usepackage{titlesec}
\usepackage{booktabs}
\usepackage{array}

\geometry{margin=2.5cm}

\definecolor{criticalcolor}{HTML}{DC3545}
\definecolor{highcolor}{HTML}{FD7E14}
\definecolor{mediumcolor}{HTML}{FFC107}
\definecolor{lowcolor}{HTML}{28A745}
\definecolor{resolvedcolor}{HTML}{6C757D}

\newcommand{\severitybox}[2]{%
  \colorbox{#1}{\makebox[2.5cm][c]{\textcolor{white}{\textbf{#2}}}}%
}

\lstset{
  basicstyle=\small\ttfamily,
  breaklines=true,
  frame=single,
  numbers=none,
  tabsize=2,
  backgroundcolor=\color{gray!10},
  extendedchars=true,
  inputencoding=utf8,
  literate={—}{---}1 {·}{$\cdot$}1 {≈}{$\approx$}1 {×}{$\times$}1 {→}{$\rightarrow$}1
}

\title{\textbf{Performance Report \#11}\\[4pt]
\large Vulkan Engine — GPU Synchronization, Avoidable Waits \&
Octree Geometry-Chunk Upload Path}
\author{Henrique Rocha}
\date{\today}

\begin{document}
\maketitle

\begin{abstract}
Where Reports \#09 and \#10 audited the shader tree, this report audits the
\emph{CPU$\leftrightarrow$GPU synchronization surface} that governs how fast the
engine can turn a freshly meshed octree node into a drawable geometry chunk on
the GPU. The scope is deliberately narrow: host-side stalls
(\texttt{vkWaitForFences}, \texttt{vkQueueWaitIdle}, \texttt{vkDeviceWaitIdle},
empty-submit readbacks), queue contention, and staging/allocation churn on the
octree chunk-streaming path (\texttt{SceneRenderer::processPendingMeshes} $\rightarrow$
\texttt{IndirectRenderer::uploadMesh*} $\rightarrow$ \texttt{runSingleTimeCommandsAsync}).
Seven issues are reported (2 high, 3 medium, 2 low). None changes rendered
output; every fix removes a host stall or allocation without altering geometry,
culling results, or feature behaviour. All recommendations respect
\texttt{AGENTS.md}: Synchronization2 is preserved, no legacy render passes or
legacy sync are introduced, feature detection is retained, and the same-queue
RADV workarounds documented in the code are kept as fallbacks.
\end{abstract}

\tableofcontents
\newpage

\section{Scope \& Methodology}

The audit followed the geometry-chunk data path end to end:

\begin{enumerate}[leftmargin=*]
  \item Background tessellation: \texttt{LocalScene::requestModel3D}
        (\texttt{utils/LocalScene.cpp:37}) meshes a node and hands a
        \texttt{Geometry} to the scene callback, which enqueues it on
        \texttt{pendingMeshQueue}.
  \item Main-thread drain: \texttt{SceneRenderer::processPendingMeshes}
        (\texttt{vulkan/renderer/SceneRenderer.cpp:1019}) pops a batch and
        uploads each chunk.
  \item GPU upload: \texttt{IndirectRenderer::uploadMeshVerticesAndIndices}
        (\texttt{vulkan/renderer/IndirectRenderer.cpp:247}) stages and copies
        vertex/index data, then defers the indirect/bounds meta-buffer write.
  \item Submission \& lifetime: \texttt{VulkanApp::runSingleTimeCommandsAsync},
        \texttt{processPendingCommandBuffers}, and the
        \texttt{deferDestroyUntilFence} recycle point.
\end{enumerate}

Every host-side wait reachable from the per-frame loop was traced to determine
whether it blocks the main (render) thread or a worker, how often it fires, and
whether the ordering it enforces could be achieved without a stall. Findings are
ordered by measured impact on chunk-streaming throughput and frame pacing.

\subsection*{Synchronization Hotspot Overview}

\begin{center}
\begin{tabular}{lcl}
\toprule
\textbf{Site} & \textbf{File:Line} & \textbf{Stall type} \\
\midrule
Per-chunk transfer publish   & \texttt{IndirectRenderer.cpp:47,361}  & \texttt{vkWaitForFences} UINT64\_MAX \\
Visible-count readback       & \texttt{IndirectRenderer.cpp:1359}    & empty submit + wait + queue idle \\
Per-chunk staging alloc      & \texttt{IndirectRenderer.cpp:341}     & \texttt{createBuffer}/free per chunk \\
Deferred fence destroy       & \texttt{IndirectRenderer.cpp:797,1398}& \texttt{vkQueueWaitIdle} in callback \\
Transfer queue routing       & \texttt{VulkanApp.cpp:1582,5851}      & single graphics queue contention \\
Legacy GPU veg. gen          & \texttt{VegetationRenderer.cpp:1488}  & \texttt{vkWaitForFences} per chunk \\
Water target reset           & \texttt{WaterRenderer.cpp:259}        & \texttt{vkDeviceWaitIdle} \\
\bottomrule
\end{tabular}
\end{center}

\newpage

\section{Issue Register}

\subsection{High Severity}

\begin{enumerate}[leftmargin=*]

\item[\textbf{[S1]}] \textbf{Per-Chunk Upload Serializes on the Previous Chunk's Transfer Fence}
  \hfill \severitybox{highcolor}{HIGH}

  \begin{tabular}{ll}
    \textbf{File:} & \texttt{vulkan/renderer/IndirectRenderer.cpp} \\
    \textbf{Line:} & 361--364 (publish), 39--62 (\texttt{publishPendingTransfer}), 47 (wait) \\
    \textbf{Category:} & Avoidable wait / octree chunk-upload throughput \\
  \end{tabular}

  \textbf{Description:}
  \texttt{processPendingMeshes} drains up to \texttt{kMaxPerFrame = 10} chunks
  per frame (\texttt{SceneRenderer.cpp:1022}, loop at 1100--1103), calling
  \texttt{uploadMesh} $\rightarrow$ \texttt{uploadMeshVerticesAndIndices} for
  each. Before recording a new staging copy, that function flushes any
  outstanding transfer:

\begin{lstlisting}
if (pendingTransfer.fence != VK_NULL_HANDLE) {
    publishPendingTransfer(app);   // IndirectRenderer.cpp:361
}
pendingTransfer.fence = app->runSingleTimeCommandsAsync(...);
\end{lstlisting}

  and \texttt{publishPendingTransfer} \emph{blocks the main thread} on the
  previous chunk's fence:

\begin{lstlisting}
if (app->resources.find((uintptr_t)pendingTransfer.fence).has_value()) {
    vkWaitForFences(dev, 1, &pendingTransfer.fence, VK_TRUE, UINT64_MAX); // :47
}
\end{lstlisting}

  Because only \emph{one} \texttt{pendingTransfer} slot exists, chunk $N+1$
  cannot begin recording until chunk $N$'s copy has fully retired on the GPU.
  The nominally ``async'' upload path (\texttt{runSingleTimeCommandsAsync}) is
  therefore run \emph{synchronously}, one round-trip per chunk. With a batch of
  10 small chunks this is 10 serialized CPU$\leftrightarrow$GPU round trips on
  the render thread every frame — the single largest limiter on how fast a
  streaming/edited octree becomes visible.

  \textbf{Recommendation:}
  Pipeline the uploads. Either (a) coalesce the whole batch into a
  \emph{single} command buffer: record all vertex/index copies for the batch,
  submit once, and publish every mesh's meta-buffer entry when that one fence
  signals; or (b) keep a small ring of in-flight \texttt{pendingTransfer} slots
  (e.g. 3, matching frames-in-flight) and only wait when the ring is full,
  reclaiming slots lazily via \texttt{pollPendingTransfers} (already called at
  \texttt{SceneRenderer.cpp:1046}). Option (a) is simplest and also amortizes
  submission overhead. Meta-buffers are append-only (documented at
  \texttt{IndirectRenderer.cpp:49--53}), so deferred publication remains correct.
  Identical GPU buffer contents; the per-chunk stall disappears.

\item[\textbf{[S2]}] \textbf{Per-Frame Visible-Count Readback: Empty Submit, Timed Wait, and a Queued \texttt{vkQueueWaitIdle}}
  \hfill \severitybox{highcolor}{HIGH}

  \begin{tabular}{ll}
    \textbf{File:} & \texttt{vulkan/renderer/IndirectRenderer.cpp}, \texttt{main.cpp} \\
    \textbf{Line:} & IndirectRenderer.cpp 1359--1416; main.cpp 1327, 1334 \\
    \textbf{Category:} & Avoidable wait / frame pacing \\
  \end{tabular}

  \textbf{Description:}
  The stats overlay reads the GPU cull count twice every frame — once for solid,
  once for water (\texttt{main.cpp:1327,1334}). Each
  \texttt{readVisibleCount} does three expensive things on the render thread:

\begin{lstlisting}
// 1) push an empty batch just to get a completion signal
vkQueueSubmit2(queue, 1, &emptySubmit, fence);        // :1380
// 2) block up to 100 ms on it
vkWaitForFences(dev, 1, &fence, VK_TRUE, 100'000'000);// :1384
// 3) UNCONDITIONALLY defer-destroy the fence, and that
//    callback forces a full-queue idle:
app->deferDestroyUntilFence(oldFence, [app, oldFence]() {
    vkQueueWaitIdle(app->getGraphicsQueue());         // :1398
    vkDestroyFence(app->getDevice(), oldFence, nullptr);
});
\end{lstlisting}

  The deferred closure runs inside \texttt{processPendingCommandBuffers}
  (\texttt{VulkanApp.cpp:1921}), so \emph{each frame enqueues two
  \texttt{vkQueueWaitIdle(graphicsQueue)} calls} that fully drain the graphics
  queue at the next safe point — defeating frames-in-flight pipelining. In
  addition, the empty-submit + 100\,ms wait is a GPU$\rightarrow$CPU sync taken
  purely to render a debug label.

  \textbf{Recommendation:}
  Make the read non-blocking and stop churning the fence. Keep a small
  ring of count buffers already indexed by \texttt{currentCullFrame}
  (\texttt{visibleCountBuffers[]} exists) and simply return the value produced
  by a frame that has \emph{already} completed (last frame's persistently-mapped
  count at \texttt{visibleCountMapped[frame]}, IndirectRenderer.cpp:1412). Drop
  the empty submit, the 100\,ms wait, and the per-call fence recreate/defer.
  This removes both the readback stall and the two per-frame queue idles; the
  displayed count lags by at most one frame, which is invisible in a stats
  overlay.

\subsection{Medium Severity}

\item[\textbf{[S3]}] \textbf{Fresh Staging Buffer Allocated (and Freed) Per Chunk}
  \hfill \severitybox{mediumcolor}{MEDIUM}

  \begin{tabular}{ll}
    \textbf{File:} & \texttt{vulkan/renderer/IndirectRenderer.cpp} \\
    \textbf{Line:} & 341--354 (alloc), 55--58 (free in \texttt{publishPendingTransfer}) \\
    \textbf{Category:} & Allocation churn / \texttt{vkAllocateMemory} minimization \\
  \end{tabular}

  \textbf{Description:}
  Each chunk upload creates a brand-new host-visible staging buffer and destroys
  it once the transfer completes:

\begin{lstlisting}
Buffer staging = app->createBuffer(stagingSize,
    VK_BUFFER_USAGE_TRANSFER_SRC_BIT,
    VK_MEMORY_PROPERTY_HOST_VISIBLE_BIT | VK_MEMORY_PROPERTY_HOST_COHERENT_BIT); // :341
...
app->resources.removeBufferVma(pendingTransfer.stagingBuffer.buffer, ...);       // :56
\end{lstlisting}

  This is a VMA allocation + map + unmap + free for every geometry chunk. The
  engine already owns a \texttt{StagingRingBuffer}
  (\texttt{vulkan/StagingRingBuffer.cpp}) with a persistently-mapped suballocator
  designed for exactly this. \texttt{AGENTS.md} explicitly calls for
  ``persistent mapped upload buffers'' and ``minimize \texttt{vkAllocateMemory}
  calls''; the current path violates both on the hottest streaming loop.

  \textbf{Recommendation:}
  Suballocate the per-chunk staging region from \texttt{StagingRingBuffer}
  (\texttt{allocate}/\texttt{release}) instead of \texttt{createBuffer}. Release
  the ring region in the transfer-completion path (or the deferred callback)
  exactly where the buffer is freed today. Combined with \textbf{S1}'s batch
  coalescing, a single ring region can back the whole batch's copies. No
  behavioural change; the per-chunk allocator traffic is eliminated.

\item[\textbf{[S4]}] \textbf{Deferred Fence Destruction Forces a Full-Queue Idle}
  \hfill \severitybox{mediumcolor}{MEDIUM}

  \begin{tabular}{ll}
    \textbf{File:} & \texttt{vulkan/renderer/IndirectRenderer.cpp}, \texttt{vulkan/VulkanApp.cpp} \\
    \textbf{Line:} & IndirectRenderer.cpp 797--802, 1395--1400; VulkanApp.cpp 1483--1489 \\
    \textbf{Category:} & Avoidable wait / queue serialization \\
  \end{tabular}

  \textbf{Description:}
  To work around RADV's validation lag on \texttt{vkDestroyFence}
  (VUID-vkDestroyFence-fence-01120), several sites recreate a fence and
  defer-destroy the old one via a closure that first drains the whole queue:

\begin{lstlisting}
app->deferDestroyUntilFence(oldFence, [app, oldFence]() {
    vkQueueWaitIdle(app->getGraphicsQueue());   // full graphics-queue stall
    vkDestroyFence(app->getDevice(), oldFence, nullptr);
});
\end{lstlisting}

  These closures execute in \texttt{processPendingCommandBuffers}
  (\texttt{VulkanApp.cpp:1921}). Every path that rotates a fence — the
  single-time ring (\texttt{VulkanApp.cpp:1486}), the visible-count fences
  (\texttt{IndirectRenderer.cpp:800}, and \textbf{S2}'s per-frame recreate at
  1398) — thus contributes a full \texttt{vkQueueWaitIdle} at a per-frame safe
  point. The idle is only there to make \emph{fence destruction} legal, not to
  order any real work.

  \textbf{Recommendation:}
  Stop destroying these fences. Maintain a per-purpose fence free-list: after a
  fence signals, \texttt{vkResetFences} it (at a point where it is provably no
  longer pending — the deferred callback already guarantees that via the
  tracking fence) and return it to the pool for reuse. Fence reset does not need
  a queue idle. Where the RADV lag genuinely requires it, destroy in a single
  batched pass guarded by one \texttt{vkGetFenceStatus} check rather than a
  \texttt{vkQueueWaitIdle} per fence. This removes the per-frame queue stalls
  while keeping validation clean.

\item[\textbf{[S5]}] \textbf{All Chunk Transfers Funnel Through the Graphics Queue; \texttt{geometryQueue} Is Acquired but Unused}
  \hfill \severitybox{mediumcolor}{MEDIUM}

  \begin{tabular}{ll}
    \textbf{File:} & \texttt{vulkan/VulkanApp.cpp} \\
    \textbf{Line:} & 1582--1595 (routing), 5851 (queue acquired) \\
    \textbf{Category:} & Queue contention / parallelism \\
  \end{tabular}

  \textbf{Description:}
  A third graphics-family queue is acquired specifically for geometry work:

\begin{lstlisting}
if (gfxRequested > 2) vkGetDeviceQueue(device, indices.graphicsFamily.value(), 2, &geometryQueue);
else geometryQueue = graphicsQueue;                    // VulkanApp.cpp:5851
\end{lstlisting}

  but the transfer entry points force everything onto \texttt{graphicsQueue}:

\begin{lstlisting}
VkFence VulkanApp::runSingleTimeCommandsAsyncOnTransfer(...) {
    return runSingleTimeCommandsAsync(fn, outSemaphore); // uses graphicsQueue
}
\end{lstlisting}

  Consequently every chunk staging copy takes \texttt{graphicsSubmitMutex} and
  competes with per-frame render submission on one queue. \texttt{geometryQueue}
  is dead. (The comments correctly note that a dedicated \emph{transfer-family}
  queue caused GPUVM faults on RADV; that is a separate hardware queue from the
  \emph{graphics-family} \texttt{geometryQueue} acquired here.)

  \textbf{Recommendation:}
  Route batched geometry-chunk copies (post \textbf{S1}) to \texttt{geometryQueue}
  when \texttt{gfxRequested > 2}, ordering the graphics consumer with a timeline
  semaphore (already the preferred primitive per \texttt{AGENTS.md}). This must
  be feature-detected: fall back to the current single-graphics-queue path when
  a distinct queue is unavailable, preserving the documented RADV safety net.
  Upload work then overlaps frame rendering instead of contending for the same
  submit mutex.

\subsection{Low Severity}

\item[\textbf{[S6]}] \textbf{Legacy GPU Vegetation Generation Still Blocks Synchronously}
  \hfill \severitybox{lowcolor}{LOW}

  \begin{tabular}{ll}
    \textbf{File:} & \texttt{vulkan/renderer/VegetationRenderer.cpp} \\
    \textbf{Line:} & 1488 \\
    \textbf{Category:} & Dead-weight wait \\
  \end{tabular}

  \textbf{Description:}
  \texttt{generateChunkInstances} dispatches the instance-gen compute shader and
  then blocks per chunk:

\begin{lstlisting}
vkWaitForFences(device, 1, &fence, VK_TRUE, UINT64_MAX); // :1488
\end{lstlisting}

  The active code path is \texttt{generateChunkInstancesCPU}
  (called at \texttt{SceneRenderer.cpp:1345}), which correctly batches all copies
  into one async submission and publishes chunks from a deferred callback
  (\texttt{VegetationRenderer.cpp:1719--1757}) — the pattern \textbf{S1} should
  adopt. The blocking GPU path appears to survive only as a fallback but still
  carries a per-chunk full stall if ever taken.

  \textbf{Recommendation:}
  Either delete the blocking GPU path or convert it to the same
  \texttt{runSingleTimeCommandsAsync} + \texttt{deferDestroyUntilFence} deferred
  publish used by the CPU path. No behaviour change on the live path.

\item[\textbf{[S7]}] \textbf{Water Target Reset Uses \texttt{vkDeviceWaitIdle}}
  \hfill \severitybox{lowcolor}{LOW}

  \begin{tabular}{ll}
    \textbf{File:} & \texttt{vulkan/renderer/WaterRenderer.cpp} \\
    \textbf{Line:} & 259 \\
    \textbf{Category:} & Avoidable device-wide stall \\
  \end{tabular}

  \textbf{Description:}
  Recreating the water offscreen targets resets the descriptor pool behind a
  full device idle:

\begin{lstlisting}
VkResult r = app->deviceWaitIdle();          // :259
if (r == VK_SUCCESS) vkResetDescriptorPool(device, waterDepthDescriptorPool, 0);
\end{lstlisting}

  This is a whole-device stall on every water-target rebuild (window resize,
  render-scale change). \texttt{AGENTS.md} restricts \texttt{vkDeviceWaitIdle} to
  shutdown and major rebuilds; a target resize is borderline but hits the entire
  device, including compute and unrelated graphics work.

  \textbf{Recommendation:}
  Prefer per-frame (frames-in-flight) descriptor pools so a rebuild only needs
  to wait on the specific in-flight fences that reference the old sets, or defer
  the pool reset through \texttt{deferDestroyUntilFence} instead of idling the
  device. If a wait is unavoidable, scope it to \texttt{queueWaitIdle} on the
  graphics queue rather than the whole device.

\end{enumerate}

\newpage

\section{Summary}

\begin{center}
\begin{tabular}{cll}
\toprule
\textbf{ID} & \textbf{Site} & \textbf{Severity} \\
\midrule
S1 & \texttt{IndirectRenderer::uploadMeshVerticesAndIndices} (per-chunk fence wait) & High   \\
S2 & \texttt{IndirectRenderer::readVisibleCount} (per-frame readback + queue idle) & High   \\
S3 & \texttt{IndirectRenderer} per-chunk staging \texttt{createBuffer} & Medium \\
S4 & \texttt{deferDestroyUntilFence} closures calling \texttt{vkQueueWaitIdle} & Medium \\
S5 & Chunk transfers on \texttt{graphicsQueue}; \texttt{geometryQueue} unused & Medium \\
S6 & \texttt{VegetationRenderer::generateChunkInstances} blocking wait & Low    \\
S7 & \texttt{WaterRenderer} descriptor-pool reset \texttt{vkDeviceWaitIdle} & Low    \\
\bottomrule
\end{tabular}
\end{center}

All seven issues are pure host-side synchronization/allocation waste: none
changes geometry, culling results, or any rendered feature. The two
highest-impact items both throttle the octree geometry-chunk pipeline —
\textbf{S1} serializes the per-frame upload batch into one-round-trip-per-chunk,
and \textbf{S2} injects a GPU readback plus two full-queue idles \emph{every
frame} for a debug counter. Fixing S1 (batch-coalesce or a small in-flight
fence ring) and S2 (return last-completed frame's mapped count, no submit/wait)
should be done first; they also make S3 (reuse \texttt{StagingRingBuffer}) and
S4 (fence free-list instead of destroy-with-idle) natural follow-ups. S5 unlocks
true upload/render overlap via the already-acquired \texttt{geometryQueue} behind
feature detection, and S6/S7 remove residual stalls on the vegetation-fallback
and resize paths. Every recommendation keeps Synchronization2, dynamic
rendering, timeline semaphores, and the documented RADV same-queue fallbacks
intact.

\end{document}
